\section{Современные методы безградиентной оптимизации}
\label{sec:modern}

\noindent\hspace{1.25cm}\emph{Раздел используется в ЛР~№8 и~№9.}

\subsection{Выбор метода оптимизации}
\label{sec:modern-choice}

Классический ГА с двоичными строками сегодня редко бывает лучшим инструментом, но выросшие из него
идеи лежат в основе методов, которые применяются в промышленности и в исследованиях искусственного
интеллекта. Общая черта этих методов~--- они используют только значения целевой функции
(<<чёрный ящик>>) и потому подходят там, где градиента нет: при подборе гиперпараметров, настройке
симуляторов, проектировании конструкций, поиске архитектур и программ. Если градиент доступен и
дёшев, градиентные методы почти всегда эффективнее. Ориентиры для выбора приведены в
таблице~\ref{tab:modern-choice}.

\begin{table}[!htb]
\caption{Выбор безградиентного метода по свойствам задачи}
\label{tab:modern-choice}
\centering\small
\begin{tabularx}{\textwidth}{>{\raggedright}p{5.2cm}>{\raggedright}p{4.3cm}>{\raggedright\arraybackslash}X}
\toprule
Задача & Метод & Библиотеки Python \\
\midrule
Непрерывная, вычисление дешёвое, до сотни переменных, много экстремумов & CMA-ES, IPOP-CMA-ES & \code{cma}, \code{CmaEsSampler} (Optuna) \\
Вычисление дорогое (обучение модели), десятки--сотни проб & байесовская оптимизация: TPE, гауссовы процессы & Optuna \\
Несколько противоречивых критериев & NSGA-II, NSGA-III, MOEA/D & pymoo, \code{NSGAIISampler} (Optuna) \\
Комбинаторная с ограничениями (маршруты, расписания) & целочисленное программирование, локальный поиск, меметический ГА & OR-Tools, HiGHS \\
Поиск формул, эвристик, программ & генетическое программирование, эволюция с LLM & PySR, FunSearch-подобные системы \\
\bottomrule
\end{tabularx}
\end{table}

\subsection{Эволюционные стратегии и CMA-ES}
\label{sec:cmaes}

\emph{Эволюционные стратегии} (ЭС) работают с вещественными векторами и основным оператором считают
мутацию $\mathbf{x}' = \mathbf{x} + \sigma\mathcal{N}(0, I)$. Главная задача ЭС~--- управлять
распределением мутаций. Простейшая (1+1)-ЭС регулирует шаг по \emph{правилу одной пятой}
И.~Рехенберга: если успешными (улучшающими) оказываются больше $1/5$ мутаций, шаг $\sigma$ увеличивают,
если меньше~--- уменьшают.

\emph{CMA-ES} (covariance matrix adaptation evolution strategy, Н.~Хансен~\cite{hansen2001, hansen2016})
адаптирует не только шаг, но и форму распределения. На каждой итерации из многомерного нормального
распределения с центром $\mathbf{m}$, шагом $\sigma$ и ковариационной матрицей $C$ порождаются $\lambda$
точек, из них отбираются $\mu = \lambda/2$ лучших, и по ним обновляются параметры распределения:
\begin{equation}
\label{eq:cma-sample}
\mathbf{x}_k = \mathbf{m} + \sigma\mathbf{y}_k, \quad \mathbf{y}_k \sim \mathcal{N}(0, C), \qquad
\mathbf{m} \leftarrow \mathbf{m} + \sigma\sum_{i=1}^{\mu} w_i\,\mathbf{y}_{i:\lambda},
\end{equation}
где $\mathbf{y}_{i:\lambda}$~--- шаг $i$-й по качеству точки, $w_1 \ge \ldots \ge w_\mu > 0$~--- веса,
$\sum w_i = 1$. Ковариационная матрица обновляется двумя слагаемыми:
\begin{equation}
\label{eq:cma-cov}
C \leftarrow (1 - c_1 - c_\mu)\,C + c_1\,\mathbf{p}_c\mathbf{p}_c^{\mathsf T}
+ c_\mu\sum_{i=1}^{\mu} w_i\,\mathbf{y}_{i:\lambda}\mathbf{y}_{i:\lambda}^{\mathsf T}.
\end{equation}
Ранг-$\mu$ слагаемое растягивает распределение в направлениях удачных шагов текущего поколения,
ранг-1 слагаемое использует \emph{путь эволюции} $\mathbf{p}_c$~--- экспоненциально сглаженную сумму
последовательных шагов среднего, которая помнит устойчивое направление движения. Шаг $\sigma$
регулируется вторым путём $\mathbf{p}_\sigma$ (\emph{кумулятивная адаптация шага}): если последовательные
шаги направлены в одну сторону, путь длиннее, чем при случайном блуждании, и шаг увеличивается, если
шаги гасят друг друга~--- уменьшается:
\begin{equation}
\label{eq:cma-sigma}
\sigma \leftarrow \sigma\exp\!\left(\frac{c_\sigma}{d_\sigma}\left(\frac{\lVert\mathbf{p}_\sigma\rVert}
{\mathrm{M}\lVert\mathcal{N}(0, I)\rVert} - 1\right)\right).
\end{equation}

В результате матрица $C$ постепенно приближает форму линий уровня функции в окрестности текущей точки
(подобно обратному гессиану в квазиньютоновских методах), и CMA-ES быстро движется вдоль узких
изогнутых оврагов (рисунок~\ref{fig:l8-cmaes} в ЛР~№8). Метод использует только ранги точек, поэтому
не зависит от монотонных преобразований $f$, и инвариантен к поворотам системы координат. Все
параметры, кроме начальных $\mathbf{m}$ и $\sigma$, имеют надёжные значения по умолчанию: $\lambda = 4 +
\lfloor 3\ln n\rfloor$, начальный шаг~--- около $0,3$ ширины области.

С популяцией по умолчанию CMA-ES ведёт себя как очень эффективный \emph{локальный} метод и на
многоэкстремальных функциях застревает в ближайшем минимуме. Стандартное решение~--- перезапуски с
увеличением популяции: \emph{IPOP-CMA-ES} после каждого останова начинает заново из случайной точки
с удвоенным $\lambda$, и большая популяция сглаживает мелкий рельеф функции.

\subsection{Байесовская оптимизация и подбор гиперпараметров}
\label{sec:bayes}

Если одно вычисление целевой функции~--- это обучение модели, бюджет исчисляется десятками проб, и
популяционные методы расходуют его слишком щедро. \emph{Байесовская оптимизация}~\cite{snoek2012} строит
по уже вычисленным точкам \emph{суррогатную} вероятностную модель функции и выбирает следующую точку,
максимизируя \emph{функцию приобретения}, которая сочетает ожидаемое качество и неопределённость. Для
гауссова процесса с прогнозом $\mu(\mathbf{x})$ и стандартным отклонением $s(\mathbf{x})$ стандартный
выбор~--- ожидаемое улучшение лучшего найденного значения $f^*$ (для минимизации):
\begin{equation}
\label{eq:ei}
EI(\mathbf{x}) = \mathrm{M}\max\bigl(0, f^* - f(\mathbf{x})\bigr) = (f^* - \mu)\,\Phi(z) + s\,\varphi(z),
\qquad z = \frac{f^* - \mu}{s},
\end{equation}
где $\Phi$ и $\varphi$~--- функция распределения и плотность стандартного нормального закона.

\emph{Древовидный парзеновский оценщик} (TPE, Дж.~Бергстра~\cite{bergstra2011})~--- алгоритм по умолчанию
в библиотеке Optuna~\cite{akiba2019}~--- делит уже проверенные точки на <<хорошие>> (лучшие $\gamma \approx
10$--$25$\,\%) и остальные, оценивает их плотности $l(\mathbf{x})$ и $g(\mathbf{x})$ и выбирает точку с
наибольшим отношением $l(\mathbf{x})/g(\mathbf{x})$, что эквивалентно максимизации ожидаемого улучшения.
TPE естественно работает с целочисленными, категориальными и условными параметрами.

Практика подбора гиперпараметров:
\begin{itemize}
\item целевая функция~--- качество на кросс-валидации, а не на обучающей выборке;
\item масштабные параметры (коэффициенты регуляризации, скорость обучения, ширина ядра) ищутся в
логарифмической шкале;
\item случайный поиск эффективнее перебора по сетке: обычно важны лишь немногие параметры, и сетка тратит
пробы на повторение одних и тех же их значений~\cite{bergstra2012};
\item безнадёжные пробы останавливают досрочно по промежуточным результатам (pruning, алгоритмы Hyperband
и ASHA);
\item лучшее значение кросс-валидации оптимистично смещено (оно выбрано как максимум среди шумных оценок),
поэтому итоговое качество измеряют на отложенной тестовой выборке или вложенной кросс-валидацией.
\end{itemize}

\subsection{Многокритериальная оптимизация}
\label{sec:moo}

В задаче с критериями $f_1, \ldots, f_M$ (например, ошибка модели и её размер, стоимость и срок
доставки) решение $\mathbf{x}$ \emph{доминирует} $\mathbf{y}$ ($\mathbf{x} \prec \mathbf{y}$), если оно
не хуже по всем критериям и строго лучше хотя бы по одному. Недоминируемые решения образуют
\emph{множество Парето}, их образы в пространстве критериев~--- \emph{фронт Парето}. Взвешенная сумма
критериев находит только выпуклую часть фронта и требует заранее знать веса, а эволюционный алгоритм
приближает весь фронт за один запуск, оставляя выбор компромисса человеку.

\emph{NSGA-II} (К.~Деб, 2002~\cite{deb2002})~--- самый распространённый многокритериальный ГА:
\begin{itemize}
\item \emph{быстрая недоминируемая сортировка} делит популяцию на фронты $F_1, F_2, \ldots$: $F_1$~---
недоминируемые решения, $F_2$~--- недоминируемые после удаления $F_1$ и т.\,д.;
\item \emph{расстояние скученности} решения внутри фронта~--- сумма нормированных расстояний между его
соседями по каждому критерию (крайним точкам~--- бесконечность); оно поощряет равномерное покрытие фронта;
\item \emph{турнир скученности}: из двух особей побеждает та, у которой меньше номер фронта, а при
равенстве~--- больше расстояние скученности;
\item элитарное сокращение ($\mu + \lambda$): из объединения родителей и потомков берутся целые фронты,
пока они помещаются, последний фронт урезается по расстоянию скученности.
\end{itemize}
Для вещественных переменных применяют \emph{имитацию двоичного кроссинговера} (SBX): потомки
$c_{1,2} = \frac12\bigl((1 \pm \beta)p_1 + (1 \mp \beta)p_2\bigr)$, где $\beta = (2u)^{1/(\eta_c + 1)}$ при
$u \le 0,5$ и $\beta = (2(1 - u))^{-1/(\eta_c + 1)}$ иначе, $u \sim U[0; 1]$, с последующим обменом
переменными между потомками с вероятностью $0,5$, и \emph{полиномиальную мутацию} с параметром
$\eta_m$. Большие $\eta$ дают потомков ближе к родителям; типичные значения $\eta_c = 15$, $\eta_m = 20$.

Качество приближения фронта измеряют \emph{гиперобъёмом} (HV)~--- площадью (объёмом) области, которую
доминирует найденное множество относительно опорной точки (больше~--- лучше), и \emph{инвертированным
расстоянием поколений} (IGD)~--- средним расстоянием от точек истинного фронта до ближайшего найденного
решения (меньше~--- лучше). Стандартные тестовые задачи~--- семейство ZDT~\cite{zitzler2000}; для
ZDT1 $f_1 = x_1$, $f_2 = g\bigl(1 - \sqrt{f_1/g}\bigr)$, $g = 1 + 9\sum_{i=2}^{n}x_i/(n - 1)$, истинный фронт
$f_2 = 1 - \sqrt{f_1}$. Для четырёх и более критериев применяют NSGA-III и MOEA/D.

\subsection{Эволюция программ с помощью языковых моделей}
\label{sec:llm-evolution}

Генетическое программирование эволюционирует программы случайными заменами поддеревьев, и большинство
таких изменений ломает программу или ничего не улучшает. Большие языковые модели (LLM), обученные на
коде, порождают синтаксически правильные и осмысленные правки, поэтому в роли оператора мутации и
скрещивания они резко повышают долю полезных потомков.

\emph{FunSearch} (DeepMind, 2023~\cite{romera2024}) организует эволюционный цикл вокруг одной функции
в программе-шаблоне:
\begin{enumerate}
\item в \emph{базе программ} (набор <<островов>> для сохранения разнообразия) хранятся версии функции
и их оценки;
\item в подсказку для LLM помещаются две--три хорошие версии, упорядоченные по возрастанию качества, и
просьба написать следующую версию;
\item сгенерированный код выполняется \emph{оценщиком} на наборе примеров, некорректные и медленные
программы отбрасываются, остальные добавляются в базу с оценкой.
\end{enumerate}
Так были найдены новые конструкции в комбинаторике (задача о cap set) и эвристики онлайн-упаковки,
превосходящие классическое правило Best Fit. Развитие подхода~--- \emph{AlphaEvolve}
(2025~\cite{novikov2025}), эволюционирующий целые файлы кода: он нашёл алгоритм умножения комплексных
матриц $4 \times 4$ за 48 скалярных умножений и эвристику планировщика задач в центрах обработки
данных Google, возвращающую в среднем 0,7\,\% вычислительных ресурсов.

Необходимые условия применения:
\begin{itemize}
\item \emph{автоматический оценщик}: качество решения должно вычисляться программой, быстро и
однозначно; без этого эволюционный поиск невозможен;
\item \emph{изоляция}: сгенерированный код считается недоверенным и выполняется в отдельном процессе с
ограничением времени, памяти и разрешённых модулей, в промышленных системах~--- в контейнере;
\item \emph{контроль переобучения}: эвристика подбирается на обучающих примерах и проверяется на
отложенных, как модель машинного обучения;
\item \emph{контроль стоимости}: число и стоимость вызовов модели фиксируются так же, как бюджет
вычислений функции.
\end{itemize}

\emph{Онлайн-упаковка в контейнеры} (online bin packing): предметы размеров $s_1, s_2, \ldots$ приходят по
одному, каждый нужно сразу и окончательно поместить в контейнер ёмкости $C$, минимизируя число
контейнеров. Простая нижняя граница $L_1 = \lceil\sum_i s_i/C\rceil$~\cite{martello1990}. Классические
правила: \emph{первый подходящий} (First Fit)~--- самый ранний контейнер, куда предмет помещается,
\emph{наилучший подходящий} (Best Fit)~--- контейнер с наименьшим остатком после укладки, \emph{наихудший
подходящий} (Worst Fit)~--- с наибольшим остатком. Эвристику удобно записать функцией приоритета
контейнеров; её и эволюционирует FunSearch.

\subsection{Реализация на Python}
\label{sec:py-modern}

Своя реализация CMA-ES на NumPy занимает около 50 строк (приложение~\ref{app:listings}); для сравнения
используют библиотеки:
\begin{lstlisting}
import cma, optuna
es = cma.CMAEvolutionStrategy(x0, 0.3 * (b - a), {"bounds": [a, b], "seed": 1})
es.optimize(f)                                   # f принимает вектор, возвращает число
study = optuna.create_study(direction="maximize", sampler=optuna.samplers.TPESampler(seed=0))
study.optimize(lambda tr: cv_score(tr.suggest_float("C", 1e-2, 1e4, log=True)), n_trials=30)

from pymoo.algorithms.moo.nsga2 import NSGA2
from pymoo.optimize import minimize
res = minimize(problem, NSGA2(pop_size=100), ("n_gen", 200), seed=1)   # res.F - фронт
\end{lstlisting}
Библиотека \code{pymoo}~\cite{blank2020} содержит также NSGA-III, MOEA/D, тестовые задачи и показатели
качества HV и IGD.
Замена сэмплера Optuna на \code{CmaEsSampler} или \code{NSGAIISampler} превращает тот же код в CMA-ES
или многокритериальный поиск. Модуль мини-FunSearch для ЛР~№9 (окружение упаковки, проверка кода,
вызов модели и эволюционный цикл) также приведён в приложении~\ref{app:listings}.
